% =====================================================================
%  RESUMEN
% =====================================================================
\chapter*{Resumen}
\addcontentsline{toc}{chapter}{Resumen}

\noindent En la Universidad Nacional de San Antonio Abad del Cusco (UNSAAC), casi ninguna experiencia de hostigamiento sexual llega a convertirse en denuncia formal. Muchas de estas situaciones quedan registradas solo como relatos escritos, y reconocer en ellos un caso de acoso exige una lectura cuidadosa, porque no siempre se describen con palabras explícitas. Esta investigación propone desarrollar un modelo de aprendizaje automático que, a partir de un relato escrito por un estudiante, determine si describe o no una situación de acoso sexual. Para ello se construirá un corpus de relatos aportados de forma voluntaria por estudiantes mayores de edad, que serán anonimizados y etiquetados por dos anotadores siguiendo una guía revisada por especialistas. El desarrollo seguirá la metodología CRISP-DM. Como modelos candidatos se entrenarán la regresión logística, la máquina de vectores de soporte y Naïve Bayes con representación TF-IDF, además del modelo BETO ajustado para el español. Todo el procesamiento se hará en equipos locales y cuidando que los datos de prueba no intervengan en el entrenamiento. El modelo elegido se comparará con una línea base léxica, basada en palabras clave, usando macro F1, precisión y exhaustividad, junto con la prueba de McNemar. La investigación es aplicada, de tipo tecnológico, con enfoque cuantitativo y diseño experimental. Como resultado se obtendrá un modelo validado y un prototipo que muestra la predicción y las palabras que más influyeron en ella, pensado como apoyo para el personal que revisa estos casos y no como sustituto de su criterio.

\vspace{0.5em}
\noindent\textit{Palabras clave}: acoso sexual, aprendizaje automático, procesamiento del lenguaje natural, clasificación de textos, educación superior.
